\chapter{Concurrency and the Failure Taxonomy}\label{ch08:concurrency}

\section{Failures are evidence too}\label{ch08:intro}

The previous chapter treated the deterministic-closure protocol as a pipeline whose happy path ends in a receipt. This chapter treats the unhappy paths, and its governing observation is simple: in an evidence kernel, a failure is not an inconvenience to be engineered away---it is a data point about the world, and the way the kernel handles it determines whether the surviving body of receipts means what it appears to mean.

Three pressures act on any system that executes subprocess workers on shared research hosts. Concurrency: multiple independent processes---test runners, notebooks, agents---launch workers on the same machine at overlapping times, and the operating system's behavior under that load is imperfect. Nondeterminism: some failures do not reproduce, and the temptation to ``just try again'' is constant. Diagnosis: when something fails, someone must be able to find out why, which pulls toward recording more; but evidence records circulate, which pulls toward recording less. The kernel's answers---an advisory lock with honest limits, an absolute no-retry rule, and failure records that are bounded and hashed rather than verbose---form a coherent position on all three, and this chapter develops each in depth. The chapter closes with the fractal bridge rule: the single, narrow, typed doorway through which legacy artifacts may touch v8 structure, included here because it is the profile's remaining fail-closed surface and because its design applies the same discipline---explicit inputs, typed failures, no silent meaning---to a semantic rather than operational problem.

Everything in this chapter, as throughout the monograph, describes a Ring-4 research candidate: the behaviors are test-evidenced under exact fixtures at the Gate~4 receipt, and nothing here is a reliability certification or an authority claim.

\section{The advisory lock discipline}\label{ch08:lock}

\subsection{The problem the lock answers}\label{ch08:lock-problem}

The lock exists because of an observed, concrete pathology, not a theoretical one. On Windows hosts, when many isolated Python interpreters initialize at the same time, a child process occasionally exits with a return code but with completely empty standard output and standard error. From the parent's perspective the worker simply vanished: no artifact, no error message, no diagnostics---an empty transport. The failure is intermittent, load-dependent, and uncorrelated with the inputs of any particular compilation.

For an ordinary build system this would be an annoyance. For an evidence procedure it is worse, because the deterministic-closure protocol runs \emph{two} workers per compilation and compares them; a launch-storm failure in either worker kills the compilation, and a storm that strikes during replay (Step 8 of Protocol~\ref{ch07:prot:closure}) kills a structural decision. The failure rate of compilations would become a function of what else the host happens to be doing---an unrecorded, uncontrolled variable modulating the production of evidence.

A process-local mutex is the obvious first move and is insufficient, for a stated reason: the processes contending for the host are not threads of one program. Independent test runners, developer shells, and research agents overlap on the same machine, and no in-process lock is visible across them. The serialization has to live somewhere all cooperating processes can see, which on a shared host means the file system.

\subsection{Why advisory}\label{ch08:why-advisory}

The chosen mechanism is a one-byte advisory file lock---\texttt{msvcrt.locking} on Windows, \texttt{fcntl.flock} elsewhere---on a file of fixed name, \texttt{cas-core-v8-compatibility-worker.lock}, in the host temporary directory, taken in exclusive non-blocking mode inside a poll loop. ``Advisory'' means the lock constrains only processes that ask: a process that never attempts the lock is not prevented by the operating system from launching whatever it likes. Why accept that?

First, because the alternative does not exist in portable form. Mandatory file locking is a niche, platform-divergent capability---effectively unavailable in any form the kernel could rely on across Windows and POSIX---and OS-level mechanisms that could truly exclude non-cooperating processes (job objects, cgroups, privileged brokers) would require exactly the sandboxing infrastructure that Boundary~\ref{ch07:bd:nonsandbox} declares out of scope for v0.1. An advisory lock is what a user-privilege, cross-platform kernel can actually have; claiming more would be false capability, and this kernel's stance on false capability fields is already on record.

Second, because the threat model is coordination, not adversaries. The empty-transport pathology arises from \emph{cooperating} software---the program's own test suites and research processes---colliding by accident. Advisory locking is precisely a coordination mechanism for cooperating parties. An adversarial process on the host could ignore the lock, but an adversarial process on the host can also do far worse things that no lock addresses (Chapter~\ref{ch07:worker-protocol} was explicit that the worker is not sandboxed); defending against accident with a mechanism proportioned to accident is the honest scope.

Third, because advisory OS locks have the one liveness property this design cannot do without: \emph{crash release}. The lock is held through an open file handle, and when the holding process dies---normally or abnormally---the operating system releases the lock. A lock-file convention based on file existence (create to acquire, delete to release) fails exactly here: a crashed holder leaves the file behind and every future launcher waits forever or implements fragile staleness heuristics. Under the OS-held advisory lock, a crashed holder stalls contenders for at most the remainder of their bounded wait, usually far less.

The kernel wraps the file lock in a second, inner discipline: a process-local thread lock acquired first, with the same bound. This is not redundancy. The file lock serializes \emph{processes}; the thread lock serializes \emph{threads within a process}, which would otherwise contend on the same file handle with less well-defined semantics, and it does so without file-system traffic for the common single-process case. The two locks share one deadline, measured on the monotonic clock---so a wall-clock adjustment mid-wait cannot extend or truncate the bound---and both failure modes surface as the same typed error, \texttt{E-COMPAT-WORKER-LOCK}, with a message distinguishing which level timed out.

Finally, the entire discipline is \emph{profile-bound}. The launch policy identifier (\texttt{cross-process-exclusive-lock-v1}), the canonical lock path, and all three timing constants are fields of the execution dependency profile, verified against the frozen module constants at profile construction and rejected on any deviation. The concurrency regime under which a receipt was produced is thus part of the receipt---a reviewer never has to guess which locking discipline was in force, and a build of the kernel with altered timing produces receipts that say so by failing validation.

\subsection{The lock file itself}\label{ch08:lockfile}

The physical lock deserves a paragraph of its own, because even here there are choices with failure modes behind them. The lock file lives at a resolved, canonical path in the host temporary directory---the one directory guaranteed writable to every cooperating process on every supported platform without configuration---under a fixed, namespaced filename. A configurable location would reintroduce the coordination problem the lock solves: two processes configured differently would lock different files and exclude nothing. The profile validator enforces canonicality, rejecting any profile whose recorded lock path is not the canonical resolution of the fixed location, so a receipt cannot claim the discipline while having used a private lock.

On first use, the acquirer opens the file in append mode---creating it if absent, never truncating it if present, since truncating a file another process holds a lock on invites platform-specific grief---and writes a single null byte if and only if the file is empty, because a byte-range lock needs at least one byte to range over. The locked region is exactly that one byte. Locking a single byte rather than the whole file keeps the operation identical across the two platform APIs, whose whole-file semantics differ, and makes the lock's extent independent of anything any process ever writes. The file's \emph{content} carries no information at all: it is not a PID file, not a timestamp, not a queue. All state lives in the OS lock table, which is exactly the state the OS reliably cleans up on process death. Every convention that stores liveness information in lock-file content eventually confronts the staleness problem---is the recorded PID still the same process?---and the kernel sidesteps the entire genre by storing nothing.

Failure handling around the lock is fail-closed in both directions. An acquirer that cannot open the lock file (permissions, exotic tempdir configuration) fails immediately with the typed lock error rather than proceeding unserialized. And---less obviously---a holder that cannot \emph{release} cleanly also raises: a release failure means the OS lock state is no longer what the discipline believes, and continuing silently would let a wedged lock manifest later as inexplicable timeouts in unrelated processes. The error surfaces at the holder, where the cause lives.

\subsection{Parameter rationale}\label{ch08:parameters}

Three numbers govern the discipline: a 90-second acquisition bound, a 10-millisecond poll interval, and a 10-millisecond post-release yield. None is magic; each is an explicit trade, and recording the reasoning is part of keeping the profile reviewable.

\emph{The 90-second bound.} The bound must be long enough that ordinary contention almost never trips it, and short enough that a wedged host fails visibly rather than hanging research pipelines indefinitely. The relevant scale is the holder's occupancy: a lock holder runs one worker under a 30-second timeout, and a full compilation runs two workers back-to-back, so a waiter behind one full compilation can legitimately wait on the order of a minute in the degenerate case where both of the holder's workers exhaust their timeouts. Ninety seconds covers that worst legitimate case with margin, while remaining small compared to the human timescales of a research session. The essential property is not the exact value but its \emph{finiteness and fail-closed character}: acquisition either succeeds within the bound or raises the typed lock error, and no code path waits forever or proceeds unlocked. An unbounded wait converts a host anomaly into a silent hang; an unlocked fallback converts it into precisely the launch storm the lock exists to prevent. Both are worse than an honest, attributable failure.

\emph{The 10-millisecond poll.} The file lock is taken in non-blocking mode inside a sleep loop rather than in blocking mode. Blocking acquisition would delegate the wait to the OS, which is cheaper, but it makes deadline enforcement platform-dependent and awkward, and it forfeits control over wakeup behavior. Polling at 10 milliseconds costs at most a hundred lock attempts per second per waiter---negligible load against a file-system lock---while bounding the latency between a release and the next acquisition at one poll interval. A much smaller interval would buy latency no compilation notices (workers run for seconds) at the price of busier spinning; a much larger one would add human-perceptible lag to serialized test suites. Ten milliseconds sits where the marginal cost of both errors is small, and, being profile-bound, it is the same everywhere evidence is produced.

\emph{The 10-millisecond post-release yield.} After releasing the file lock, the holder sleeps ten milliseconds before returning. This is the most unusual of the three and exists because of a scheduling asymmetry: the thread that just released is warm---running, cache-resident, and (in the common test-suite pattern) about to loop around and want the lock again---whereas a waiter in another process must be scheduled before its next poll can succeed. Without the yield, a tight local loop can reacquire the lock indefinitely ahead of a cross-process waiter, not through any queue-jumping in the OS but simply because it asks more often from a running state. The yield hands the scheduler an explicit opportunity to run someone else first.

\subsection{What is honestly not claimed}\label{ch08:no-fairness}

It must be stated plainly, because the receipt-consuming layers of the kernel rely on statements being exact: this discipline claims \emph{mutual exclusion among cooperating launchers within a bounded wait}, and nothing else. It does not claim fairness---waiters may acquire in any order, and the OS lock APIs used here maintain no ordered queue the kernel could appeal to. It does not claim starvation-freedom---an unlucky waiter can lose every race for the full 90 seconds and then fail, and the post-release yield is a mitigation that improves the odds, not a mechanism that bounds them. It does not claim anything about non-cooperating processes, which are invisible to an advisory lock by definition.

Why belabor non-claims? Because of how properties propagate. A downstream engineer who reads ``serialized'' as ``fair'' will eventually build something---a latency budget, a liveness argument, a test---on the fairness they inferred, and it will fail rarely, unreproducibly, and far from its cause. The program's documented failure history is largely a history of inferred properties hardening into assumed ones. The kernel's defense is to make its concurrency contract exactly as strong as its mechanism, in writing, in the profile: \emph{bounded exclusive launch, unfair, advisory, fail-closed}. A future version that wants fairness must build a mechanism that has it, and change the policy identifier when it does.

\begin{ruleenv}[Concurrency contract]\label{ch08:rule:contract}
Cooperating worker launches are mutually exclusive under the profile-bound advisory lock; acquisition is bounded at 90 seconds with 10-millisecond polling and fails closed; a releasing holder yields 10 milliseconds. No fairness, starvation-freedom, or exclusion of non-cooperating processes is claimed, and no admitted receipt asserts any.
\end{ruleenv}

\section{The no-retry principle}\label{ch08:no-retry}

The protocol never retries a failed worker. Not on lock timeout, not on empty transport, not on nonzero exit, not on malformed output, not on the nondeterminism failures of Section~\ref{ch08:disagreement}. A failed worker fails its compilation, and the failure propagates, typed, to the caller. This section develops the argument for that rule in full, because it is the point where the kernel's evidence discipline conflicts most sharply with ordinary engineering instinct, and because the rule is load-bearing for the meaning of every receipt the kernel emits.

The instinct says: transient failures are noise; retrying filters noise; filtered results are cleaner. The instinct is correct for \emph{operations} and wrong for \emph{evidence}, and the difference is worth stating carefully.

Consider what a compatibility receipt is for. Chapter~\ref{ch07:worker-protocol} defined the deterministic-closure property: within the recorded perimeter, compilation is a deterministic function, tested by running two workers per compile and replaying before every decision. The population of receipts is therefore not merely a set of useful outputs---it is the \emph{measurement record} of that determinism claim. Every compilation is a trial; agreement is a confirming instance; disagreement or worker failure is an anomaly that bears on whether the closure property actually holds on real hosts.

Now introduce retry. A retry policy re-runs failed trials and keeps the successes. The surviving record then no longer estimates the behavior of the system; it estimates the behavior of the system \emph{conditioned on eventual success}. Failures whose causes are real but intermittent---a dependency file being rewritten during runs, an interpreter behaving differently under memory pressure, a compiler bug that manifests one run in fifty, a host process interfering with workers---are exactly the failures retry masks best, because they are exactly the failures that pass on the second attempt. The receipts look clean; the closure property looks confirmed; and the anomalies that would have falsified it have been quietly deleted by the harness itself. This is selection bias in its textbook form---the analyst sees only the survivors---except that here the selection is applied automatically, invisibly, and before any analyst is involved. No field in a retried receipt says ``this succeeded on attempt three''; and if such a field existed and were honest, the receipt would no longer be the receipt of a deterministic procedure at all.

The statistical analogy is optional stopping: re-running an experiment until it yields the desired outcome and reporting only the final run invalidates the reported result, even though every individual run was honestly conducted. Retry-until-success is optional stopping implemented in infrastructure. The kernel's refusal is the refusal to let its own harness perform the maneuver.

There is a second, independent argument from \emph{identity}. A receipt binds one execution: these bytes, this profile, this artifact. Under retry, what would the receipt describe? The successful attempt only---with the failed attempts either discarded (evidence destroyed) or appended (in which case the record describes a composite procedure, ``run until success within $n$ tries,'' whose determinism is a different and weaker property than the one the kernel tests). Either way the crisp statement ``a compatibility execution is a deterministic function of its recorded inputs'' is lost. The two-worker agreement check makes this concrete: its entire evidential force comes from the fact that both workers are unconditional. Two workers that agree \emph{after the harness discarded disagreeing pairs} would demonstrate nothing except the harness's persistence.

Third, there is an argument from \emph{attribution}. The launch-storm pathology of Section~\ref{ch08:lock-problem} was discovered, characterized, and mitigated because it produced visible failures. A retry layer above the worker would have absorbed it years earlier at small per-run cost---and the cost would have compounded invisibly: an unserialized, retrying system fails more, retries more, loads the host more, and fails more again, with no signal reaching anyone. Fail-closed systems surface their pathologies while the pathologies are cheap; fail-retry systems surface them as capacity mysteries later. The advisory lock exists \emph{because} the no-retry rule made the underlying failure impossible to ignore---the two disciplines are complements, the lock reducing the failure rate honestly (by removing a cause) rather than cosmetically (by hiding occurrences).

The objection writes itself: is a compilation that fails because the host was busy not simply wasted work? It is---and the waste is priced correctly. The caller sees a typed failure with a transport record; the caller may absolutely decide to run a new compilation, and nothing forbids a human or an orchestration layer from doing so. The difference between that and in-protocol retry is where the decision is visible. A caller-level rerun is a new, independent compilation, producing its own complete evidence, with the failed attempt's record still existing wherever the caller's logs put it; the decision to rerun was made \emph{above} the evidence layer by an accountable party. An in-protocol retry is made \emph{inside} the evidence layer by the harness, and leaves no trace in the surviving record. The kernel's position is not that reruns are forbidden; it is that the evidence layer itself must be memoryless and unconditional, so that whatever selection happens, happens where it can be seen.

A final clarification distinguishes this rule from the retry discipline of distributed-systems practice, because the surface similarity misleads. In a distributed system, retrying an idempotent request against a replicated service is sound engineering: the request's \emph{outcome} is the object of interest, the service's transient unavailability is genuinely irrelevant to that outcome, and the retry semantics are part of the declared contract---callers know results are at-least-once outcomes and design accordingly. None of the three conditions holds here. The compatibility protocol's object of interest is not the compiled artifact alone but the \emph{execution as an event}: the receipt's value lies in describing what one run of the closed procedure did, which is exactly the thing retry redefines. Worker failures are not presumed irrelevant---the determinism claim makes every failure potentially relevant, and presuming otherwise is the bias. And no consumer of a receipt is in a position to ``design around'' hidden retries, because the receipt format asserts single-execution semantics; a retrying implementation would make every receipt a small lie. Retry is not wrong in general; it is wrong for artifacts whose meaning is ``this is what happened,'' and the kernel's receipts have no other meaning.

\begin{invariant}[Unconditional workers]\label{ch08:inv:no-retry}
Every worker launch admitted by the profile is unconditional: its outcome, success or failure, propagates to the caller as the outcome of its compilation. No admitted code path launches a worker because a previous worker failed.
\end{invariant}

\section{Worker transport failure records}\label{ch08:transport}

When a worker fails at the transport level---wrong exit status, undecodable output, stderr traffic, protocol-shape violation---the kernel must record something, and the question is what. Its answer is deliberately spare: the failure record embeds exactly the worker's exit code, the byte counts of standard output and standard error, and the SHA-256 digests of both streams. Five numbers and no content.

The exclusions are the design. Stream \emph{content} is excluded for three reasons that compound. Bound: a misbehaving worker can emit arbitrarily much; a failure record that embeds content has no size bound, and failure records travel---into exceptions, logs, test output, evidence packets---so an unbounded record is an unbounded liability everywhere downstream. Sensitivity: stderr from a failing process is precisely the kind of text that leaks incidental host detail---user paths, environment fragments, library versions---and Section~\ref{ch08:disagreement} will make the general principle explicit: evidence records must not exfiltrate host contents. Trust: the content of a failed worker's streams is, by construction, output of software in an unknown state; embedding it verbatim in records that conforming tools will parse and display extends the blast radius of whatever produced it.

What the five fields preserve is exactly what diagnosis and integrity need. The exit code classifies the failure at the OS level. The byte counts distinguish the signature cases at a glance---the launch-storm pathology is \emph{exit code with zero bytes on both streams}, instantly recognizable---and separate ``said nothing'' from ``said something malformed.'' The digests do the quiet heavy lifting: they make the failure record checkable against any copy of the streams that does exist elsewhere (a caller that captured them can prove they are the same bytes the record describes), they make identical failures identical as records (two occurrences of the same misbehavior produce matching digests, so recurrence is detectable across time and hosts without content comparison), and they commit the record to a specific observed behavior without reproducing it. A transport record is thus a content-addressed \emph{reference} to the failure, in the same sense that the rest of the kernel's evidence is content-addressed reference to artifacts \cite{reprobuilds}.

The fail-closed casework around the transport record was described in Chapter~\ref{ch07:worker-protocol} and is completed here from the failure side. An empty response is never treated as an empty success: no parseable protocol object means no worker claim at all, and the compilation fails with the transport record attached. A malformed response---undecodable bytes, duplicate keys, wrong shape, wrong protocol identifier---likewise fails closed; the parser's strictness (Rule~\ref{ch07:rule:exact-shapes}) means there is no lenient path along which a damaged response could be partially believed. Contradictory responses fail on the contradiction itself: a worker that reports success but exits nonzero, or that produces a perfect success document while also writing to stderr, is a worker whose self-description is inconsistent, and inconsistency is grounds for rejection, not adjudication. In every case the error carries its typed code---\texttt{E-COMPAT-WORKER-EXIT}, \texttt{E-COMPAT-WORKER-PROTOCOL}, \texttt{E-COMPAT-WORKER-STDERR}, \texttt{E-COMPAT-WORKER-START}, \texttt{E-COMPAT-WORKER-LOCK}---so the failure taxonomy remains enumerable, and, per Invariant~\ref{ch08:inv:no-retry}, none of them triggers a second attempt.

\begin{ruleenv}[Transport minimality]\label{ch08:rule:transport}
A worker transport failure is recorded as exit code, stdout and stderr byte counts, and SHA-256 digests of both streams---nothing else. Stream content never enters an evidence record. Empty and malformed responses fail closed and are never retried.
\end{ruleenv}

\section{Dependency-profile disagreement}\label{ch08:disagreement}

The subtlest failure in the taxonomy is the one where nothing failed. Both workers of a compilation launch cleanly, execute the verified compiler bytes, and return well-formed success responses---and their execution dependency profiles differ. Perhaps a Python module file changed between the two launches; perhaps the two processes resolved different search paths; perhaps something in the interpreter's initialization is not as deterministic as assumed. Whatever the cause, the premise of the receipt---that ``the execution environment'' is a single well-defined thing the receipt describes---has just been falsified by experiment, and the compilation must die: \texttt{E-COMPAT-NONDETERMINISTIC-PROFILE}.

(The sibling failure, in which the profiles agree but the compiled \emph{artifacts} differ, is \texttt{E-COMPAT-NONDETERMINISTIC-RESULT}; it needs no special record design, because the artifacts' digests already identify the disagreement. The profile case is harder precisely because profiles describe the host.)

The error record for a profile disagreement contains: both profile digests, and a diagnostic listing \emph{at most 32} keyed difference paths, with a boolean flag stating whether the list was truncated. The comparison runs over a transformed view of each profile in which every environment variable's value is replaced by its hash, and every dependency's file path is likewise replaced by its hash. Three design decisions---the bound, the keying, the hashing---each carry an argument.

\emph{Why the bound.} A difference enumeration between two structured documents has no natural size limit: if the profiles diverge early and broadly---say, entirely different module inventories---a naive diff enumerates hundreds of paths, and the diagnostic dwarfs the evidence it annotates. The failure record inherits every liability of unbounded records from Section~\ref{ch08:transport}, with one addition special to diagnostics: a huge diff is also \emph{useless}, because a reviewer confronting three hundred difference paths learns less than one confronting the first thirty and a flag saying ``and more.'' Thirty-two paths suffice to localize any plausible real-world divergence (a changed file, a differing path set, a drifted environment entry produce a handful of paths each); the truncation flag preserves honesty about what was cut; and the bound makes the worst-case record size a constant of the profile rather than a function of how strange the host was. The traversal itself is deterministic---dictionary keys visited in sorted order, lists by index with an explicit length-difference marker---so the \emph{same} 32 paths are reported for the same disagreement, and two occurrences of one underlying cause produce comparable diagnostics.

\emph{Why keyed paths.} The paths are expressed against keyed collections---environment entries addressed as \texttt{name=PYTHONUTF8}, dependencies as \texttt{module\_name=json}---rather than positional indices. Positional addressing is fragile under exactly the perturbation being diagnosed: if one profile contains an extra module, every subsequent index shifts, and a positional diff reports dozens of spurious differences downstream of one real insertion. Keyed addressing pins each reported difference to the entity it concerns, so an inserted module is one path, not an avalanche.

\emph{Why hashes for values.} Environment values and dependency paths are the two fields in a profile that describe the host rather than the computation, and they are exactly the fields excluded, as values, from the diagnostic. The reasons parallel Section~\ref{ch08:transport} but sharpen. An environment value on a research host can contain usernames, machine names, directory conventions, and---in misconfigured but common cases---credentials; a dependency path embeds the local file-system layout. A profile-disagreement diagnostic is generated precisely when something unexpected is true of the host, which is the worst possible moment to copy host detail into a record that will be attached to an exception, pasted into an issue, or bundled into an evidence packet. Hashing preserves the two facts a diagnostic needs---\emph{whether} the values differ (different hashes) and \emph{which entity} differs (the keyed path)---while making the record safe to circulate. A reviewer with legitimate access to the host can recompute the hashes and identify the value; a reader of the record alone cannot recover it. The asymmetry is the point: diagnosis is possible for those already inside the trust perimeter, and the record extends no one's reach beyond it \cite{anderson}.

\begin{ruleenv}[Bounded, keyed, hashed disagreement]\label{ch08:rule:disagreement}
A dependency-profile disagreement between two successful workers records both profile digests and at most 32 deterministic, keyed difference paths with an explicit truncation flag; environment values and dependency paths appear in the diagnostic only as hashes. The compilation fails closed and is not retried.
\end{ruleenv}

\begin{remark}[The diagnostic view is itself canonical]\label{ch08:rem:diag-canonical}
The transformed comparison view---environment values and dependency paths replaced by hashes, collections re-keyed by name---is constructed by the same canonical machinery as everything else in the kernel: deterministic key order, canonical text, one serialization route \cite{jcs}. This matters for a reason easy to miss. The diagnostic exists to be \emph{compared}: against the diagnostic of a past disagreement, against a teammate's, against one generated on another host. If the view's construction were ad hoc---iteration-order dependent, locale-sensitive, or version-drifting---two reports of one underlying divergence would differ textually, and the recurrence detection that makes intermittent environmental faults findable would silently degrade. Applying the evidence encoding to the failure diagnostics closes the loop: even the kernel's descriptions of its own anomalies are stable, hashable artifacts.
\end{remark}

It bears repeating what this failure's existence does for the kernel's epistemics. A system that ran one worker would never observe profile disagreement; its receipts would assert environmental determinism as an assumption. The two-worker design converts the assumption into a per-compilation experiment, and the disagreement record is the experiment's falsification report---bounded and hashed because even falsification reports are evidence, and evidence discipline does not relax when the news is bad.

\section{The fractal bridge rule}\label{ch08:fractal}

The chapter's final subject sits at the boundary between the legacy profile and the v8 kernel proper. Legacy CAS artifacts contain \emph{fractal digit sequences}---ordered sequences of digits 0 through 9 whose interpretation, in the legacy system, is positional and implicit. Chapter~\ref{ch07:worker-protocol} noted that the compatibility profile flags any artifact containing a fractal literal with the diagnostic \texttt{UNTYPED\_FRACTAL\_REQUIRES\_BRIDGE} and supplies no interpretation of its own. The bridge is the one admitted mechanism by which such a sequence may be given v8-structural form, and its design is a compact restatement of the kernel's whole philosophy.

The bridge is a function taking two inputs: the digit sequence, and a \emph{caller-supplied catalog} mapping digits to typed finite arrows of the TKS v8 reference kernel---the finite semigroupoid of typed arrows developed by the companion candidate. Its output, when it succeeds, is a \texttt{SemigroupoidPath}: the typed composition of the arrows selected by the digits, in sequence order. Every departure from that happy path fails closed with a typed error: a sequence that is not an ordered, replayable sequence of items (strings and byte strings are explicitly rejected, since Python would happily iterate them as if they were sequences of digits); a digit outside 0--9, or a boolean masquerading as an integer; a digit for which the catalog supplies no mapping (\texttt{E-COMPAT-FRACTAL-UNMAPPED}); a catalog entry that is not a typed finite arrow; and---the semantically load-bearing case---a digit-adjacent pair of arrows whose target and source types do not meet, which the path constructor rejects and the bridge surfaces as \texttt{E-COMPAT-FRACTAL-TYPE}. There is no default catalog, no partial result, and no untyped fallback.

Two design choices deserve their arguments.

\emph{Why the catalog is caller-supplied.} The kernel could have shipped a canonical digit-to-arrow mapping---after all, the legacy system ``knows'' what its digits mean. It does not, because that knowledge is precisely the kind of cross-namespace correspondence that the semantic-identity chapter forbids treating as identity. A claim of the form ``legacy digit 7 corresponds to v8 arrow $f$'' is a translation hypothesis between ontologies; under the kernel's rules such a claim must be represented as a typed link carrying its own source and review state, never baked silently into infrastructure. A built-in catalog would be a silent, unreviewable, kernel-endorsed translation of the entire legacy fractal vocabulary---the largest semantic overclaim the kernel could possibly make, installed at the exact chokepoint every fractal must pass through. Requiring the caller to supply the catalog relocates the semantic commitment to where it belongs: with a party who can be asked ``on what authority does 7 map to $f$?'', whose answer is subject to review, and whose choice is visible in every use. The kernel's own contribution is limited to what it can actually verify: that the supplied mapping is well-typed and that the resulting composition is admissible in the finite reference kernel.

\emph{What composition proves, and what it does not.} When the bridge returns a path, exactly one thing has been established: the digit sequence, \emph{under the caller's catalog}, selects arrows that compose---each arrow's target matches its successor's source, so the sequence denotes a well-formed path in the typed finite semigroupoid. This is a real, checkable, non-trivial structural fact; sequences fail it, and the failure is informative. But it is a fact about the catalog and the type structure, not about the legacy digits. The bridge proves type-correct composition; it never proves, suggests, or records that the composition \emph{means} what the legacy system meant by the digits. A different catalog over the same sequence yields a different path with equal structural standing. Any claim that a particular catalog is the \emph{right} translation of legacy fractal semantics is an empirical and historical claim about the legacy system, to be established---if ever---by the review machinery, and no output of the bridge is evidence for it.

\begin{ruleenv}[Bridge scope]\label{ch08:rule:bridge}
Legacy fractal digit sequences acquire v8 structure only through an explicit caller-supplied catalog of typed finite arrows, composed as a semigroupoid path. Missing digits, unmapped digits, untyped catalog entries, and source/target mismatches fail closed. A successful bridge establishes type-correct composition under the supplied catalog and nothing else; in particular, it establishes no meaning of any legacy digit.
\end{ruleenv}

\begin{example}[A bridge run, end to end]\label{ch08:ex:bridge}
Suppose a legacy artifact contains the fractal literal with digit sequence $(2, 7, 7)$, and a caller supplies a catalog mapping digit $2$ to a typed arrow $f : A \to B$ and digit $7$ to $g : B \to B$, both drawn from the finite reference kernel. The bridge validates the sequence (ordered, integer digits in range, no booleans), resolves each digit through the catalog, and asks the path constructor to compose $(f, g, g)$. The types meet---$f$ lands in $B$, where both copies of $g$ begin and end---so the constructor returns the path, and the caller now holds a typed object witnessing that \emph{under this catalog} the sequence denotes a composable chain $A \to B \to B \to B$.

Now vary the run. If the catalog omits digit $7$, the bridge fails with \texttt{E-COMPAT-FRACTAL-UNMAPPED} at the first $7$, naming the digit; nothing partial is returned. If the caller instead maps $7$ to an arrow $h : C \to C$, the constructor rejects the pair $(f, h)$---target $B$, source $C$---and the bridge surfaces \texttt{E-COMPAT-FRACTAL-TYPE} with the constructor's reason. If the caller passes the string \texttt{"277"} as the sequence, the bridge rejects it before looking at any digit, because admitting string iteration would let a text field silently masquerade as a digit sequence. And if the caller's catalog maps every digit to arrows on a single object, every sequence composes---which illustrates the scope rule from the other side: composability under a trivializing catalog is a fact about the catalog, and the bridge's output never claims otherwise.

What the example shows is the division of epistemic labor. The caller owns the semantic commitment (this catalog is my hypothesis about digit meaning); the kernel owns the structural verification (under your hypothesis, this sequence does or does not compose, with a typed reason); and the boundary between the two is visible in every artifact the bridge produces or refuses to produce.
\end{example}

One implementation footnote completes the picture: the bridge imports the v8 category types lazily, at call time, so the legacy compilation path itself never loads the v8 kernel. The dependency direction of the companion-candidate architecture---Core depends on no companion's semantics---is thus preserved even at module-import granularity: a host that compiles legacy sources and never bridges a fractal never touches v8 category code at all, and the loaded-module inventory of its receipts shows exactly that.

\section{The failure codes, collected}\label{ch08:codes}

For reference, Table~\ref{ch08:tab:codes} collects the principal typed failure codes of the compatibility profile, grouped by the stage of the protocol that raises them. The grouping itself is informative: reading down the table is reading the deterministic-closure protocol as a sequence of ways to fail closed, and the reviewer confronted with any code can locate, from the table alone, which guarantee was being enforced when the run died. The list reflects the v0.1 reference implementation at the Gate~4 receipt; it is a closed set in the sense that conforming code raises no compatibility failure outside it, and any extension of the set is a reviewable change to the profile.

\begin{longtable}{@{}lll@{}}
\caption{Principal typed failure codes of the \texttt{legacy\_only} compatibility profile.}\label{ch08:tab:codes}\\
\toprule
Code & Stage & Marks \\
\midrule
\endfirsthead
\toprule
Code & Stage & Marks \\
\midrule
\endhead
\bottomrule
\endfoot
\texttt{E-COMPAT-SOURCE-HASH} & intake & source bytes contradict declared digest \\
\texttt{E-COMPAT-VERSION} & intake & declared version outside frozen policy \\
\texttt{E-COMPAT-PROFILE} & intake & declared syntax profile not admitted \\
\texttt{E-COMPAT-ENVELOPE} & intake & source envelope wrong type or shape \\
\texttt{E-COMPAT-POLICY} & policy & frozen policy constants incoherent \\
\texttt{E-COMPAT-COMPILER-MISSING} & policy & compiler bytes absent or unreadable \\
\texttt{E-COMPAT-COMPILER-IDENTITY} & policy & retained bytes contradict identity \\
\texttt{E-COMPAT-COMPILER-RAW-HASH} & policy & raw digest outside frozen variants \\
\texttt{E-COMPAT-COMPILER-CONTENT-HASH} & policy & LF content digest not frozen value \\
\texttt{E-COMPAT-COMPILER-SOURCE} & policy & source authority not canonical \\
\texttt{E-COMPAT-WORKER-MISSING} & launch & packaged worker script absent \\
\texttt{E-COMPAT-WORKER-LOCK} & launch & bounded lock acquisition or release failed \\
\texttt{E-COMPAT-WORKER-START} & launch & worker process could not execute \\
\texttt{E-COMPAT-WORKER-EXIT} & transport & exit status contradicts response \\
\texttt{E-COMPAT-WORKER-STDERR} & transport & worker emitted stderr traffic \\
\texttt{E-COMPAT-WORKER-PROTOCOL} & transport & malformed or wrongly shaped response \\
\texttt{E-COMPAT-RUNTIME} & binding & runtime identity not derivable or drifted \\
\texttt{E-COMPAT-ISOLATION} & binding & isolation flags, lock policy, or paths wrong \\
\texttt{E-COMPAT-ENVIRONMENT} & binding & environment not the sanitized profile \\
\texttt{E-COMPAT-DEPENDENCY-PROFILE} & binding & module inventory malformed or untyped \\
\texttt{E-COMPAT-DEPENDENCY-ESCAPE} & binding & file module outside admitted roots \\
\texttt{E-COMPAT-DEPENDENCY-RACE} & binding & bound file bytes changed since capture \\
\texttt{E-COMPAT-ENGINE} & artifact & engine other than Python-only profile \\
\texttt{E-COMPAT-UNBOUND-IMPORT} & artifact & \texttt{use} import or external CAS import \\
\texttt{E-COMPAT-ARTIFACT-SHAPE} & artifact & inadmissible value shape in output \\
\texttt{E-COMPAT-ARTIFACT-TEXT} & artifact & noncanonical text in output \\
\texttt{E-COMPAT-ARTIFACT-NUMBER} & artifact & nonfinite numeric in output \\
\texttt{E-COMPAT-ARTIFACT-HASH} & artifact & frozen artifact contradicts digest \\
\texttt{E-COMPAT-DIAGNOSTICS} & artifact & diagnostic tuple missing or untyped \\
\texttt{E-COMPAT-NONDETERMINISTIC-RESULT} & closure & paired workers' artifacts differ \\
\texttt{E-COMPAT-NONDETERMINISTIC-PROFILE} & closure & paired workers' profiles differ \\
\texttt{E-COMPAT-REPLAY} & closure & replay input not an exact compilation \\
\texttt{E-COMPAT-REPLAY-MISMATCH} & closure & fresh replay differs from record \\
\texttt{E-COMPAT-BOUNDARY} & boundary & artifact not Ring 4 / OPEN / NOT LICENSED \\
\texttt{E-COMPAT-V8-TRANSLATION} & boundary & record carries a v8 translation \\
\texttt{E-COMPAT-FRACTAL-SEQUENCE} & bridge & sequence not ordered and replayable \\
\texttt{E-COMPAT-FRACTAL-DIGIT} & bridge & digit outside 0--9 or boolean \\
\texttt{E-COMPAT-FRACTAL-CATALOG} & bridge & catalog missing or entry untyped \\
\texttt{E-COMPAT-FRACTAL-UNMAPPED} & bridge & digit has no catalog mapping \\
\texttt{E-COMPAT-FRACTAL-TYPE} & bridge & arrows do not compose \\
\end{longtable}

Two observations about the table close the survey. First, the codes partition cleanly by protocol stage with essentially no overlap, which is a property of the staged design rather than of the naming: each guarantee is enforced at one place, so each violation has one home. Second, the \emph{closure} rows are the profile's empirical heart---they are the codes that can fire only because the kernel does strictly more work than correctness requires (a second worker, a fresh replay), and each firing is a fact about the world that a cheaper design would never have observed.

\section{The taxonomy as a designed surface}\label{ch08:taxonomy}

Stepping back, the failure taxonomy of the compatibility profile is best understood as a designed surface with three properties, each argued somewhere above and collected here.

\emph{Enumerable.} Every failure is a typed code from a closed set---source and identity failures (Chapter~\ref{ch07:worker-protocol}), lock and transport failures (Sections~\ref{ch08:lock}--\ref{ch08:transport}), nondeterminism failures (Section~\ref{ch08:disagreement}), bridge failures (Section~\ref{ch08:fractal})---never a bare string invented at the failure site. Enumerability is what lets failure \emph{rates by kind} be measured across the receipt population, which the no-retry rule guarantees is an unbiased record of trials.

\emph{Bounded.} Failure records have constant-bounded size: five transport fields, two digests plus at most 32 paths, one typed code with canonical text. No failure, however exotic the host state that produced it, can generate an evidence artifact whose size or content depends on that state in unbounded ways.

\emph{Non-exfiltrating.} Failure records reference host particulars---stream content, environment values, file paths---only through digests. The records are safe to attach, forward, and archive by default, and diagnosis deepens only for parties who already hold the referenced bytes.

A failure surface with these properties does quiet work for the whole program. It means the unhappy paths of the kernel are as reviewable as the happy paths; it means an anomaly on one host in one month is recognizably \emph{the same anomaly} on another host in another month; and it means that the growth of the receipt corpus---successes and failures alike---is growth in usable evidence rather than in liability. The next chapter ascends one level, to the machinery that consumes these receipts and compilations: the runtime decision model, where the kernel's refusal to convert structure into authority is made total.
